% 20260708(1)_Assingment_ComputationalMathematics_A
\documentclass[dvipdfmx]{jsreport} 
\usepackage{soupreamble}

\title{計算数学特論A 第5回レポート課題}
\author{M1 6012 川村聡一郎}
\date{}


\begin{document}
\maketitle

\begin{tcolorbox}
\begin{enumerate}
	\item 独立集合問題(判定版)が多項式時間で解ければ頂点被覆問題(判定版)も多項式時間で解けることを示せ。
	\item 頂点被覆問題(判定版)が多項式時間で解ければ、独立集合問題(判定版)も多項式時間で解けることを示せ。
	\item 頂点被覆問題(探索版)は入力が木グラフのとき多項式時間で解けることを示せ。
\end{enumerate}
\end{tcolorbox}

\begin{souanswer} %答案はここに書く
	答案はここに書く
\end{souanswer}












\end{document}